\section{Retaining an exact growing exponential core}
\label{sec:exponential-core-perturbation}

An inverse--logarithmic expansion of the Lambert function does not resolve
the much smaller correction caused by adding $x^2$ to $xe^x$. The operation
\wl{AsymptoticExponentialCoreInverse} retains the Lambert inverse exactly
and expands those corrections in a separate small exponential scale.
Its zero sector is an exact inverse, and each retained coefficient is a
complete function of that inverse.

\subsection{The admitted core and source charts}

Use a positive coordinate $v\to+\infty$ and consider the exact model
\begin{equation}
 F(v)=F_0(v)+R(v),\qquad
 F_0(v)=y_0+a v^b e^{c v^p},
 \qquad c,p>0,\quad a\ne0,
 \label{eq:exponential-core-forward}
\end{equation}
where $b$ is a fixed real number and $R$ is a finite real power--log
expression $\sum_\alpha v^\alpha P_\alpha(\log v)$.
Every such perturbation, including its finitely many derivatives, is
smaller than the growing core sufficiently near the endpoint.

At an infinite original source endpoint, the reconstruction is
$x=H(v)=s+\sigma v$ with $\sigma\in\{1,-1\}$. The source translation
$s$ can be supplied explicitly or recognized from affine factors in the
core. At a finite endpoint, use $x=H(v)=x_0+\sigma/v$.
Thus the same operation covers both $xe^x+x^2$ at infinity and
$e^{1/x}+x$ at $0^+$. Define $r_*=1$ for the first reconstruction and
$r_*=-1$ for the second.

Let $Y=(y-y_0)/a>0$. The selected exact core inverse is
\begin{equation}
 v_0=
 \begin{cases}
 (\log Y/c)^{1/p},&b=0,\\[0.3em]
 \displaystyle
 \left(\frac{b}{cp}
 W_\nu\!\left(\frac{cp}{b}Y^{p/b}\right)\right)^{1/p},&b\ne0,
 \end{cases}
 \quad
 \nu=\begin{cases}0,&b>0,\\-1,&b<0.\end{cases}
 \label{eq:exponential-core-exact-inverse}
\end{equation}
These are the real branches with $v_0\to+\infty$.
For $b<0$, the argument is negative and the lower branch selects the
increasing large-$v$ part of the core. The target domain records the real
Lambert argument condition, positivity of the real-power base, and
$b+cpv_0^p>0$. In all cases the exact inverse in the original source is
$\phi(y)=H(v_0(y))$.

An optional supplied \wl{"CoreInverse"} is retained only after a bounded
symbolic comparison establishes equality with this selected exact
expression. Additional terms inside the stated core are rejected if they
would make \eqref{eq:exponential-core-exact-inverse} merely an approximation.
They can instead be placed in the separately declared finite perturbation.

\subsection{Complete coefficients without expanding the core inverse}

Put
\begin{equation}
 D(v)=a v^{b-1}(b+cpv^p),\qquad
 F_0'(v)=e^{cv^p}D(v),\qquad E(v)=e^{-cv^p}.
 \label{eq:exponential-core-derivative}
\end{equation}
Introduce the marker $\lambda$ in $F_0(v)+\lambda R(v)=y$.
The exact-core Lagrange formula of Section~\ref{sec:core-perturbations}
shows that its degree-$n$ correction has the form
$C_n(v_0)E(v_0)^n$. The coefficient $C_n$ can be computed using only
power--logarithmic and rational operations on $v$:
\begin{align}
 Q_{n,1}(v)&=\frac{H'(v)R(v)^n}{D(v)},
 \label{eq:exponential-core-coefficient-initial}\\
 Q_{n,j+1}(v)&=
 \frac{Q_{n,j}'(v)-jcpv^{p-1}Q_{n,j}(v)}{D(v)},
 \qquad 1\le j<n,\label{eq:exponential-core-coefficient-recursion}\\
 C_n(v)&=\frac{(-1)^n}{n!}Q_{n,n}(v).
 \label{eq:exponential-core-coefficient-final}
\end{align}
Indeed, differentiating $Q(v)E(v)^j$ gives
$(Q'-jcpv^{p-1}Q)E^j$, and division by $F_0'$ adds one factor of $E$.
The implementation therefore avoids intermediate derivatives of
\wl{ProductLog} and avoids mixing the exact exponential factor into each
coefficient calculation.

The exact core identity gives a second, often more convenient expression
for the small scale:
\begin{equation}
 E_0=e^{-cv_0^p}=\frac{v_0^b}{Y}.
 \label{eq:exponential-core-target-sector}
\end{equation}
The result records both expressions and uses the right-hand side in the
retained formula. It is positive on the selected branch. Through inclusive
sector degree $N$, the approximation is
\begin{equation}
 \phi(y)+\sum_{n=1}^N C_n(v_0(y))E_0(y)^n.
 \label{eq:exponential-core-retained-sectors}
\end{equation}
Here one complete marker degree is also one exponential degree, because
the forward perturbation contains no additional exponentials.

\begin{example}[The inverse of $xe^x+x^2$]
Let $v_0=W_0(y)$ and $E_0=v_0/y$. The first coefficients are
\begin{align}
 C_1(v)&=-\frac{v^2}{v+1},
 \label{eq:xexpx-core-first-coefficient}\\
 C_2(v)&=\frac{v^3(-v^2+2v+4)}{2(v+1)^3}.
 \label{eq:xexpx-core-second-coefficient}
\end{align}
Thus the first correction to the exact core inverse is
\[
 (xe^x+x^2)^{-1}(y)
 =W_0(y)-\frac{W_0(y)^3}{(W_0(y)+1)y}+\cdots.
\]
This is much smaller than the error of any fixed inverse--logarithmic
truncation of $W_0(y)$. Retaining $W_0$ exactly makes that correction
meaningful without requiring a second, deeper core approximation.
\end{example}

\subsection{A uniform complete-tail theorem}

For the finite perturbation in \eqref{eq:exponential-core-forward}, choose
\begin{equation}
\begin{gathered}
 d=\max(0,\max_\alpha(\alpha-b)),\qquad
 k=\max_\alpha\deg P_\alpha,\\
 B(v)=v^d(1+|\log v|)^k,\qquad \chi(v)=B(v)e^{-cv^p}.
\end{gathered}
 \label{eq:exponential-core-envelope}
\end{equation}
The zero perturbation is treated separately and leaves the exact inverse
unchanged.

\begin{theorem}[Complete tail around a growing exact core]
\label{thm:exponential-core-complete-tail}
For fixed data in \eqref{eq:exponential-core-forward}, there exist positive
constants $C,K$ and a source threshold such that the selected real inverse
has a convergent expansion \eqref{eq:exponential-core-retained-sectors}
with its complete omitted tail bounded by
\begin{equation}
 |\operatorname{Tail}_N|
 \le C v_0^{r_*-p}
 \frac{(K\chi(v_0))^{N+1}}{1-K\chi(v_0)},
 \qquad K\chi(v_0)<1.
 \label{eq:exponential-core-geometric-tail}
\end{equation}
In particular, for fixed $N$ the remainder class is
\begin{equation}
 \OO\!\left(E_0^{N+1}
 v_0^{r_*-p+(N+1)d}(1+|\log v_0|)^{(N+1)k}\right).
 \label{eq:exponential-core-asymptotic-tail}
\end{equation}
The bound covers the whole omitted series independently of cancellation
in the first omitted coefficient.
\end{theorem}
\begin{proof}
Set
\[
 v=v_0(1+\epsilon V),\qquad\epsilon=v_0^{-p}.
\]
The displacement scale $v_0^{1-p}$ keeps the change in the exponential
phase bounded. Dividing the forward equation by $a v_0^b e^{cv_0^p}$
gives a core difference
\begin{equation}
 (1+\epsilon V)^b
 \exp\left(c\frac{(1+\epsilon V)^p-1}{\epsilon}\right)-1.
 \label{eq:exponential-core-normalized-equation}
\end{equation}
It extends analytically to $\epsilon=0$, and its derivative in $V$ at
$(\epsilon,V)=(0,0)$ is $cp\ne0$.

The normalized perturbation is
$E_0R(v_0(1+\epsilon V))/(a v_0^b)$.
After extraction of its powers of $v_0$ and $\log v_0$, each of its finite
coefficient factors is analytic and uniformly bounded near
$\epsilon=0,V=0$, with $1/\log v_0$ as an additional small parameter.
Its absolute envelope is a constant times $B(v_0)E_0$.
Replace $E_0$ temporarily by a complex variable $Z$ and rescale
$\zeta=B(v_0)Z$. The remaining normalized coefficient parameters range
over a bounded finite polydisc. Since the derivative $cp$ at $Z=0$ is
independent of them, the analytic implicit function theorem gives a
uniform solution near $V=0$ for $|\zeta|$ smaller than a fixed radius.
Equivalently, it is analytic on $|Z|<1/(K B(v_0))$.

For the observable, subtract its exact core value. The nonconstant part is
\[
 H(v_0(1+\epsilon V))-H(v_0)
 =\sigma v_0^{r_*-p}
 \frac{(1+\epsilon V)^{r_*}-1}{\epsilon}.
\]
The quotient is analytic near $(\epsilon,V)=(0,0)$ and uniformly bounded
on a smaller polydisc. Cauchy estimates bound its $Z^n$ coefficient by
$C(KB)^n$, with the displayed factor $v_0^{r_*-p}$.
The coefficient identity above identifies these coefficients with the
exact-core Lagrange coefficients. Summing the geometric tail and setting
$Z=E_0$ proves \eqref{eq:exponential-core-geometric-tail}.
Finally $\chi(v_0)\to0$, and the eventual derivative of the real forward
function has the core's nonzero sign, so this solution is the selected
real inverse near the endpoint.
\end{proof}

For $xe^x+x^2$, one can take $d=1,k=0,p=1,r_*=1$.
The complete tail after sector $N$ is consequently
$\OO((v_0E_0)^{N+1})$.
For $e^{1/x}+x$ at $0^+$, the reconstruction has $r_*=-1$,
and the first correction is
\begin{equation}
 (e^{1/x}+x)^{-1}(y)
 =\frac1{\log y}+\frac1{y\log^3y}
 +\OO\left(\frac1{y^2\log^2y}\right).
 \label{eq:finite-endpoint-exponential-core-example}
\end{equation}
The last bound is conservative but follows directly from
\eqref{eq:exponential-core-asymptotic-tail} with $d=0$.

\subsection{An unknown input error remains a separate first sector}

Suppose the actual perturbation includes a real continuously differentiable
unknown term $A(v)$ satisfying
\begin{equation}
 A(v)=\OO(v^{-\rho}(1+|\log v|)^\ell),\qquad
 A'(v)=\OO(v^{-\rho-1}(1+|\log v|)^\ell).
 \label{eq:exponential-core-input-error}
\end{equation}
Here $\rho$ may be any fixed real number: the growing exponential core
dominates every such finite power scale. Its derivative is asymptotic to
$acp\,v^{b+p-1}e^{cv^p}$. The mean value transport argument, with the
derivative of $H$, gives an additional inverse error
\begin{equation}
 \OO\!\left(E_0v_0^{r_*-b-p-\rho}
                  (1+|\log v_0|)^\ell\right).
 \label{eq:exponential-core-input-transport}
\end{equation}
The matching derivative assumption ensures eventual monotonicity and
permits comparison on the same local inverse branch. The model inverse
and the actual inverse are both asymptotic to the exact core coordinate,
so the scale can be expressed in $v_0$.

Equation~\eqref{eq:exponential-core-input-transport} has exponential degree
one, even when the finite model has been expanded through a much larger
sector depth. The implementation therefore stores it separately and adds
it to the model-tail remainder. It does not relabel it as a higher sector
or claim that additional exact model coefficients overcome unknown input
data. The retained terms beyond this error still describe the explicitly
specified finite model.

\subsection{Implemented behavior and verification}

\begin{lstlisting}[language=Mathematica]
AsymptoticExponentialCoreInverse[
  x Exp[x], x^2, {x, Infinity}, {y, 2}]

AsymptoticExponentialCoreInverse[
  7 + 3 (x - 2) Exp[x - 2], 5 (x - 2)^2,
  {x, Infinity}, {y, 1}]

AsymptoticExponentialCoreInverse[
  Exp[1/(2 - x)], 2 - x, {x, 2}, {y, 1},
  Direction -> "FromBelow"]
\end{lstlisting}

The result identifies scale \wl{"ExactExponentialCoreSectors"}.
It records the exact inverse and branch certificate, the divergent source
coordinate and reconstruction, the complete sector coefficients, the first
omitted coefficient, both representations of $E_0$, and the separate model
and input-error remainders. Its \wl{"MajorantContract"} is an asymptotic
existence theorem, with \wl{"NumericCertificate"} set to \wl{False}; the
constants and uniform threshold are not computed numerical bounds.

Independent regression checks insert
$v+\sum_{n=1}^N C_n(v)Z^n$ into
$X e^{X-v}-v+ZX^2=0$ for the primary example and perform ordinary Taylor
expansion in $Z$. Further checks use original forward targets generated
from known source points, compare successive sector depths at high
precision, and exercise affine shifts and scales, negative source infinity,
the lower real Lambert branch, nonunit phase powers and finite endpoint
reconstruction. Rejection cases cover a decaying core, an additional term
incorrectly placed in the exact core, an extra exponential rank in the
perturbation, an unverified inverse, and unproved coefficient reality.

The admitted phase is a single growing monomial $cv^p$. General phase
polynomials would need an appropriate exact inverse supplied with its own
branch contract, while perturbations such as $e^{v/2}$ need more than the
single exponential sector used here. The operation does not replace those
missing contracts with a finite inverse--logarithmic core approximation.
